
\subsection{Performance Evaluation}

\paragraph{Important note on units.}
Our evaluation is performed in a \emph{comparative simulation} built from the provided logs. Reported waiting-time metrics are therefore best interpreted as \emph{simulation-time units} induced by the assumed travel/door parameters. The primary decision question in this study is the \emph{relative ordering and consistency} across parking strategies; absolute magnitudes can be re-scaled once building-specific timing parameters are calibrated in deployment. Accordingly, we emphasize distributional and tail metrics (e.g., P95 and long-wait rate) that remain informative under uniform re-scaling of timing parameters.
Similarly, tail-event rates depend on the chosen long-wait threshold; we therefore report P95/P99 and stress-test comparisons to avoid over-interpreting zero rates under a single nominal threshold.


\paragraph{Simulator consistency and calibration.}
To ensure credibility of simulation-based evaluation, all strategies are tested on the \emph{same} call stream with an identical simulator:
the same travel time per floor, door time, and 5-minute decision epoch are applied throughout. These parameters are selected within
realistic ranges for multi-floor elevator service and are used only to establish a consistent time scale; absolute magnitudes can be
re-calibrated once building-specific timing specifications are available.

We perform two practical sanity checks. First, under static baselines (\emph{last-stop} and \emph{lobby-return}), the simulator produces
non-degenerate service dynamics (no instantaneous service; queues form only under sufficiently high demand), and waiting-time distributions
are qualitatively consistent with observed operational variability. Second, we perturb timing parameters within a moderate band (e.g.,
$\pm 20\%$) and find that the qualitative ordering among strategies is preserved, indicating that conclusions are driven by policy structure
(mode conditioning and demand-weighted parking) rather than fragile calibration.


\paragraph{Task 1 (Forecasting).}
Figure~\ref{fig:task1-pred} shows that the one-step-ahead model tracks short-term fluctuations of hall-call arrivals at 5-minute resolution.
Table~\ref{tab:task1-metrics} reports MAE and RMSE on a chronological split, supporting the forecast as an actionable early-warning signal for upcoming surges and regime transitions.

\paragraph{Task 2 (Mode discovery and online classification).}
Table~\ref{tab:task2-cluster-profile} summarizes cluster feature profiles, which we translate into manager-friendly traffic modes (e.g., \emph{Idle/Low}, \emph{Up-Peak}, \emph{Down-Peak}, \emph{Inter-floor}).
Figure~\ref{fig:task2-heatmap} demonstrates clear time-of-day regularity, indicating that the discovered modes correspond to stable building usage patterns rather than random noise.
Because mode boundaries can overlap in real operations, we do not expect perfect cluster separation; instead, we prioritize interpretability and downstream control stability.
Figure~\ref{fig:task2-heatmap} also indicates consistent daily mode blocks, which motivates mode-conditioned parking policies that react when the building shifts between these blocks.


\paragraph{Task 3 (Mode-aware dynamic parking).}
Table~\ref{tab:task3-strategy} compares three parking policies: \emph{last-stop} (stay), \emph{lobby-return}, and our \emph{dynamic mode-aware} policy.
Across the same call stream, the dynamic policy improves service consistency by positioning idle elevators closer to likely call origins under the current regime, thereby reducing expected response distance and lowering the frequency of extreme waits.
The benefit is most pronounced during peak periods and short transitions (e.g., pre-peak buildup) where static rules either over-concentrate cars at the lobby or fail to cover dispersed demand.
To make the operational trade-off explicit, we also report repositioning cost via empty travel, enabling service--wear trade-offs in deployment.



\paragraph{Interpretable baseline benchmark and deployment fallback.}
Beyond static baselines, we implement the interpretable Route B policy in Section~\ref{subsec:baseline-fallback} as an additional
end-to-end reference. Its forecasting and state rules are fully transparent, enabling quick diagnosis when performance shifts.
In deployment, it also supports a low-risk rollback: if mode classification confidence degrades or key signals are missing, the controller
can revert to Route B without interrupting in-service elevators. For clarity, we focus the main comparison table on the two standard static baselines; Route B details and implementation notes are provided in Appendix~\ref{app:routeb}.

\subsection{Stress Tests and Robustness}
\label{subsec:stress}

Beyond average performance, elevator service quality is dominated by \emph{tail risk} (P95 and, where reported, P99 waiting times, as well as the long-wait rate).
We therefore conduct three stress tests that deliberately violate typical daily patterns, including abrupt demand spikes, mixed-direction
traffic, and regime confusion during transitions. Each test is evaluated under identical simulator parameters to ensure a fair comparison. 
Operationally, we generate stress cases by injecting additional calls into selected windows and/or reweighting directional composition while keeping the same simulator and dispatch constraints.


Overall, Route~A remains consistently strong on tail metrics, while Route~B provides a stable, low-variance fallback with controlled
repositioning cost. Appendix~\ref{app:extra} reports the full stress-test table; we highlight two operational observations:
\begin{enumerate}[label=(\alph*)]
    \item the largest gains appear during regime transitions where proactive parking shortens the initial response distance, 
    \item when mode confidence degrades, Route~B's zone allocator avoids oscillatory relocations and preserves acceptable tail performance.
\end{enumerate}

\subsection{Robustness and Sensitivity}

We verify that performance improvements are not due to fragile tuning by varying key design parameters:
\begin{itemize}[leftmargin=1.2em,itemsep=0.25em,topsep=0.25em]
  \item \textbf{Number of clusters $K$.} Varying $K$ in a small range (e.g., $4$--$8$) preserves the daily mode structure and does not change the qualitative strategy ordering, indicating that the parking optimizer is not overly sensitive to mode granularity.
  \item \textbf{Decision frequency.} Recomputing targets every 5 minutes versus less frequent updates (e.g., 10 minutes) yields a smooth trade-off: less frequent updates reduce empty repositioning but react more slowly to regime transitions. This supports deployability because operators can choose a frequency consistent with wear-and-tear constraints.
  \item \textbf{Feature smoothing window.} Using 10--20 minute rolling windows stabilizes online classification. Dominant regimes remain stable, while borderline slices may shift between similar clusters; because repositioning applies only to idle elevators, these boundary fluctuations have limited operational impact.
  \item \textbf{Timing-parameter perturbation.} Scaling travel/door parameters within a reasonable band (e.g., $\pm 20\%$) preserves the relative comparisons across strategies, reinforcing that conclusions are structural rather than tied to a single calibration.
\end{itemize}

Overall, the pipeline's gains arise from its architecture (mode conditioning + demand-weighted facility location) rather than narrowly tuned hyperparameters.
This supports implementation readiness because operators can adjust $K$ and decision frequency to meet wear-and-tear constraints without changing the qualitative ranking.
